%------------------------------------------------------------------------------%
\documentclass[fleqn,aspectratio=169,12pt]{beamer}

\usepackage{gaal}

\title{Cônicas Rotacionadas}

%------------------------------------------------------------------------------%
\begin{document}

\addtitle

%------------------------------------------------------------------------------%
\begin{frame}{Cônicas}
  Equação geral
  \[
    ax^2 + 2bxy + cy^2 + dx + ey + f = 0
  \]

  \pause
  Forma matricial
  \[
    v^tAv \;+\; B^tv \;+\; f = 0
  \]
  \pause
  com
  \[
    \pause
    v = \begin{bmatrix} x \\ y \end{bmatrix}
    \qquad
    \pause
    A = \begin{bmatrix}
      a & b \\
      b & c
    \end{bmatrix}
    \qquad
    \pause
    B = \begin{bmatrix} d \\ e \end{bmatrix}
    \qquad
    \pause
    f \in \R
  \]
\end{frame}

%------------------------------------------------------------------------------%
\begin{frame}{Identificação de Uma Cônica}

  \emph{Forma canônica}

  \pause
  \qquad cônica centrada na origem e alinhada com os eixos

  \pause
  \bigskip
  \emph{Cônica transladada}

  \pause
  \qquad é colocada na forma canônica pela transformação de coordenadas
  \[
    \pause
    \qquad
    X = x - h
    \qquad
    Y = y - k
  \]
  \pause
  \qquad pode ser necessário completar quadrados

  \pause
  \bigskip
  \emph{Se a cônica estiver rotacionada?}
\end{frame}

%------------------------------------------------------------------------------%
\framedgraphic*{Translação e Rotação}{transformacao_rotacao-1}{}
\framedgraphic*{Translação e Rotação}{transformacao_rotacao-2}{}
\framedgraphic*{Translação e Rotação}{transformacao_rotacao-3}{}
\framedgraphic*{Translação e Rotação}{transformacao_rotacao-4}{}
\framedgraphic*{Translação e Rotação}{transformacao_rotacao-5}{}

%------------------------------------------------------------------------------%
\begin{frame}{Rotação}
  A rotação da cônica está relacionada ao termo misto \(2bxy\)

  \pause
  \bigskip
  Se \(b=0\) temos uma cônica alinhada com os eixos

  \pause
  \bigskip
  Precisamos diagonalizar a matriz $A$ que é simétrica
\end{frame}

%------------------------------------------------------------------------------%
\begin{frame}{Diagonalização da Parte Quadrática}
  Como \(A\) é \emph{simétrica}

  \pause
  As matrizes

  \pause
  $D$ \emph{diagonal},
  \pause
  com os autovalores de $A$
  \[
    \pause
    D = \diag(\lambda_1, \lambda_2)
  \]

  \pause
  $P$ \emph{ortogonal},
  \pause
  com os autovetores de $A$
  \[
    \pause
    P = \left[ \, v_1\;, v_2\, \right]
  \]
  \pause
  Satisfazem
  \[
    \pause
    A = P D P^t
    \pause
    \qquad\text{ou}\qquad
    P^t A P = D
  \]
\end{frame}

%------------------------------------------------------------------------------%
\begin{frame}{Mudança de Variáveis}
  \onslide<+->{Definimos novas variáveis \(w=\vetor{\mu}{\nu}\)}
  \[
    \onslide<+->{w = P^t v}
    \onslide<+->{\qquad\text{ou}\qquad v = Pw}
  \]
  \onslide<+->{Então}
  \begin{align*}
    \onslide<+->{v^t A v}
    \onslide<+->{& = v^t P D P^t v \\}
    \onslide<+->{& = (P^tv)^t D P^t v \\}
    \onslide<+->{& = w^t D w}
  \end{align*}
  \onslide<+->{Nas novas variáveis, $w$,}
  \onslide<+->{a matriz não possui termo misto}
\end{frame}

%------------------------------------------------------------------------------%
\begin{frame}{Transformação da Equação Matricial}
\begin{columns}
\column{0.46\textwidth}
  \onslide<+->{Aplicando a transformação \(v = Pw\)}
  \begin{align*}
    \onslide<+->{v^tAv \;+\; B^tv \;+\; f               & = 0 \\}
    \onslide<+->{(Pw)^tA(Pw) \;+\; B^t(Pw) \;+\; f      & = 0 \\}
    \onslide<+->{w^t (P^t A P) w \;+\; (B^tP) w \;+\; f & = 0 \\}
    \onslide<+->{\emph{w^t D w \;+\; E^t w \;+\; f}     & = \emph{0}}
  \end{align*}
\column{0.4\textwidth}
  \emph{\[
    \onslide<+->{D \;=\; P^t A P}
    \onslide<+->{  \;=\; \diag(\lambda_1, \lambda_2)}
  \]}
  \emph{\[
    \onslide<+->{E \;=\; PB}
    \onslide<+->{  \;=\; \left[ \, v_1\; v_2\, \right] \begin{bmatrix} d \\ e \end{bmatrix}}
  \]}
\end{columns}
\end{frame}

%------------------------------------------------------------------------------%
%------------------------------------------------------------------------------%
\begin{question}
  Considerando a equação
  \[
    3x^2 + 4xy + 3y^2 - 15 = 0
  \]
  \vspace{-\baselineskip}
  \begin{enumerate}
    \item escreva a equação na forma matricial
    \item diagonalize a matriz da parte quadrática
    \item faça a mudança de variáveis
    \item identifique a cônica
  \end{enumerate}
\end{question}

%------------------------------------------------------------------------------%
\begin{resolution}{Equação Matricial}
  Equação
  \[
    3x^2 + 4xy + 3y^2 - 15 = 0
  \]
  Equação na forma matricial
  \[
    v^tAv + B^tv + f = 0
  \]
  \[
    v = \begin{bmatrix} x \\ y \end{bmatrix}
    \qquad
    A = \begin{bmatrix}
      3 & 2 \\
      2 & 3
    \end{bmatrix}
    \qquad
    B = \begin{bmatrix} 0 \\ 0 \end{bmatrix}
    \qquad
    f = -15
  \]
\end{resolution}

%------------------------------------------------------------------------------%
\begin{resolution}{Diagonalizando $A$ -- Polinômio Característico}
  Polinômio característico
  \begin{align*}
    p(\lambda)
    & = \det(A-\lambda I)
    \\
    & = \begin{vmatrix}
      3-\lambda & 2 \\
      2 & 3-\lambda
    \end{vmatrix}
    \\
    & = (3-\lambda)^2 - 4
    \\
    & = 9 - 6\lambda + \lambda^2 - 4
    \\
    & = \lambda^2 - 6\lambda + 5
  \end{align*}
\end{resolution}

%------------------------------------------------------------------------------%
\begin{resolution}{Diagonalizando $A$ -- Autovalores}
  Equação característica
  \[
    \lambda^2 - 6\lambda + 5 = 0
  \]
  \[
    (\lambda-5)(\lambda-1)=0
  \]
  Autovalores
  \[
    \lambda_1=5
    \qquad
    \lambda_2=1
  \]
\end{resolution}

%------------------------------------------------------------------------------%
\begin{resolution}{Autovetores Associados a $\lambda_1=5$}
  \begin{align*}
    (A-5I)v_1 & = \vec{0} \\
    \begin{bmatrix}
      -2 & 2 \\
      2 & -2
    \end{bmatrix}
    \vetor{x}{y}
    & = \vetor{0}{0}
  \end{align*}
  Portanto
  \[
    -2x+2y=0
    \qquad
    y=x
  \]
  Escolhendo \(x=1\)
  \[
    v_1=\vetor{1}{1}
  \]
\end{resolution}

%------------------------------------------------------------------------------%
\begin{resolution}{Autovetores Associados a $\lambda_2=1$}
  \begin{align*}
    (A-I)v_2 & = \vec{0} \\
    \begin{bmatrix}
      2 & 2 \\
      2 & 2
    \end{bmatrix}
    \vetor{x}{y}
    & = \vetor{0}{0}
  \end{align*}
  Portanto
  \[
    2x+2y=0
    \qquad
    y=-x
  \]
  Escolhendo \(x=1\)
  \[
    v_2=\vetor{1}{-1}
  \]
\end{resolution}

%------------------------------------------------------------------------------%
\begin{resolution}{Matriz Ortogonal}
  Normalizando os autovetores
  \[
    u_1=\frac{1}{\sqrt{2}}\vetor{1}{1}
    \qquad
    u_2=\frac{1}{\sqrt{2}}\vetor{1}{-1}
  \]
  Matriz ortogonal
  \[
    P=\left[u_1\;u_2\right]
    =\frac{1}{\sqrt{2}}
    \begin{bmatrix}
      1 & 1 \\
      1 & -1
    \end{bmatrix}
  \]
  Matriz diagonal
  \[
    D=P^tAP
    =\begin{bmatrix}
      5 & 0 \\
      0 & 1
    \end{bmatrix}
  \]
\end{resolution}

%------------------------------------------------------------------------------%
\begin{resolution}{Mudança de Variáveis}
  Mudança de variáveis
  \[
    v=Pw
    \qquad
    w=\vetor{\mu}{\nu}
  \]
  \[
    \begin{bmatrix}
      x \\ y
    \end{bmatrix}
    =
    \frac{1}{\sqrt{2}}
    \begin{bmatrix}
      1 & 1 \\
      1 & -1
    \end{bmatrix}
    \begin{bmatrix}
      \mu \\ \nu
    \end{bmatrix}
  \]
  Portanto
  \[
    x=\frac{\mu+\nu}{\sqrt{2}}
    \qquad
    y=\frac{\mu-\nu}{\sqrt{2}}
  \]
\end{resolution}

%------------------------------------------------------------------------------%
\begin{resolution}{Equação Transformada}
  Como \(B=\vec{0}\), temos \(E=P^tB=\vec{0}\).

  \begin{align*}
    w^tDw+E^tw+f & = 0
    \\
    \begin{bmatrix}\mu&\nu\end{bmatrix}
    \begin{bmatrix}
      5&0\\
      0&1
    \end{bmatrix}
    \begin{bmatrix}\mu\\\nu\end{bmatrix}
    -15 & = 0
    \\
    5\mu^2+\nu^2 & = 15
  \end{align*}
\end{resolution}

%------------------------------------------------------------------------------%
\begin{resolution}{Identificação}
  Dividindo a equação por \(15\)
  \[
    \frac{\mu^2}{3}+\frac{\nu^2}{15}=1
  \]
  A cônica é uma elipse com centro na origem.

  Semieixos
  \[
    a=\sqrt{15}
    \qquad
    b=\sqrt{3}
  \]

  O eixo maior está na direção de \(\nu\)
\end{resolution}

%------------------------------------------------------------------------------%
\begin{resolution}{Esboço}
  \centering
  \includegraphics{I-08-Rotacao-ex-1}
\end{resolution}
%------------------------------------------------------------------------------%




\newcommand{\mat}[1]{\begin{bmatrix}#1\end{bmatrix}}
\newcommand{\vet}[1]{\begin{bmatrix}#1\end{bmatrix}}

\begin{question}
  \begin{block}{}
    Considerando a equação
    \[
      x^2+4xy+y^2-2x-2y-3=0
    \]
  \end{block}

  \begin{enumerate}
    \setlength{\itemsep}{0.45cm}
    \item escreva a equação na forma matricial
    \item diagonalize matriz da parte quadrática
    \item faça a mudança de variáveis

      (apresente a transformação e transforme a equação)
    \item identifique a cônica
  \end{enumerate}
\end{question}

\begin{resolution}{Equação Matricial}
  \textbf{Equação}
  \[
    x^2+4xy+y^2-2x-2y-3=0
  \]
  \[
    ax^2+2bxy+cy^2+dx+ey+f=0
  \]

  \textbf{Equação na forma matricial}
  \[
    v^tAv+B^tv+f=0
  \]

  \[
    v=\mat{x\\y}
    \qquad
    A=\mat{a&b\\b&c}
     =\mat{1&2\\2&1}
    \qquad
    B=\mat{d\\e}
     =\mat{-2\\-2}
    \qquad
    f=-3
  \]
\end{resolution}

\begin{resolution}{Diagonalizando A -- Polinômio Característico}
  \textbf{Polinômio característico}
  \[
    p(\lambda)=\det(A-\lambda I)
  \]
  \[
    =\det\left(
      \mat{1&2\\2&1}
      -\lambda\mat{1&0\\0&1}
    \right)
  \]
  \[
    =\begin{vmatrix}
      1-\lambda&2\\
      2&1-\lambda
    \end{vmatrix}
  \]
  \[
    =(1-\lambda)^2-4
  \]
  \[
    =1-2\lambda+\lambda^2-4
  \]
  \[
    =\lambda^2-2\lambda-3
  \]
\end{resolution}

\begin{resolution}{Diagonalizando A -- Equação Característica}
  \textbf{Equação característica}
  \[
    p(\lambda)=\lambda^2-2\lambda-3=0
  \]
  \[
    \Delta=b^2-4ac=(-2)^2-4\times1(-3)=16
    \qquad
    \sqrt{\Delta}=4
  \]
  \[
    \lambda_1=\frac{-b+\sqrt{\Delta}}{2a}
    =\frac{2+4}{2}=3
  \]
  \[
    \lambda_2=\frac{-b-\sqrt{\Delta}}{2a}
    =\frac{2-4}{2}=-1
  \]

  \textbf{Autovalores}
  \[
    \lambda_1=3
    \qquad
    \lambda_2=-1
  \]
\end{resolution}

\begin{resolution}{Autovetores Associados ao Autovalor $\lambda_1=3$}
  \begin{columns}[T]
    \begin{column}{0.58\textwidth}
      \[
        (A-3I)v_1=\vec{0}
      \]
      \[
        \left(
        \mat{1&2\\2&1}
        -3\mat{1&0\\0&1}
        \right)
        \mat{x\\y}
        =
        \mat{0\\0}
      \]
      \[
        \mat{-2&2\\2&-2}
        \mat{x\\y}
        =
        \mat{0\\0}
      \]
    \end{column}
    \begin{column}{0.37\textwidth}
      \[
        -x+y=0
      \]
      \[
        y=x
      \]
      Escolhendo $x=1$
      \[
        v_1=\mat{1\\1}
      \]
    \end{column}
  \end{columns}
\end{resolution}

\begin{resolution}{Autovetores Associados ao Autovalor $\lambda_2=-1$}
  \begin{columns}[T]
    \begin{column}{0.58\textwidth}
      \[
        (A-(-1)I)v_2=\vec{0}
      \]
      \[
        \left(
        \mat{1&2\\2&1}
        +\mat{1&0\\0&1}
        \right)
        \mat{x\\y}
        =
        \mat{0\\0}
      \]
      \[
        \mat{2&2\\2&2}
        \mat{x\\y}
        =
        \mat{0\\0}
      \]
    \end{column}
    \begin{column}{0.37\textwidth}
      \[
        2x+2y=0
      \]
      \[
        y=-x
      \]
      Escolhendo $x=1$
      \[
        v_2=\mat{1\\-1}
      \]
    \end{column}
  \end{columns}
\end{resolution}

\begin{resolution}{Matriz Ortogonal}
  \textbf{Normalizando o autovetor $v_1$}
  \[
    \|v_1\|=\sqrt{1^2+1^2}=\sqrt{2}
    \qquad
    u_1=\frac{v_1}{\|v_1\|}
    =\frac{1}{\sqrt{2}}\mat{1\\1}
  \]

  \textbf{Normalizando o autovetor $v_2$}
  \[
    \|v_2\|=\sqrt{1^2+(-1)^2}=\sqrt{2}
    \qquad
    u_2=\frac{v_2}{\|v_2\|}
    =\frac{1}{\sqrt{2}}\mat{1\\-1}
  \]

  \textbf{Matriz ortogonal $P$ e matriz diagonal $D$}
  \[
    P=[u_1\ u_2]
    =\frac{1}{\sqrt{2}}
      \mat{1&1\\1&-1}
    \qquad
    D=\operatorname{diag}(\lambda_1,\lambda_2)
    =\mat{3&0\\0&-1}
  \]
\end{resolution}

\begin{resolution}{Mudança de Variáveis}
  \textbf{Mudança de variáveis $v=Pw$}
  \[
    \mat{x\\y}
    =
    \frac{1}{\sqrt{2}}
    \mat{1&1\\1&-1}
    \mat{\mu\\\nu}
    =
    \frac{1}{\sqrt{2}}
    \mat{\mu+\nu\\\mu-\nu}
    =
    \mat{\dfrac{\mu+\nu}{\sqrt{2}}\\[2mm]
         \dfrac{\mu-\nu}{\sqrt{2}}}
  \]
  \[
    x=\frac{\mu+\nu}{\sqrt{2}}
    \qquad
    y=\frac{\mu-\nu}{\sqrt{2}}
  \]
\end{resolution}

\begin{resolution}{Transformando as Matrizes}
  \[
    D=P^tAP=\operatorname{diag}(\lambda_1,\lambda_2)
    =\mat{3&0\\0&-1}
  \]

  \[
    E=P^tB
    =\frac{1}{\sqrt{2}}
      \mat{1&1\\1&-1}
      \mat{-2\\-2}
    =\frac{1}{\sqrt{2}}
      \mat{-4\\0}
    =\mat{-2\sqrt{2}\\0}
  \]
\end{resolution}

\begin{resolution}{Equação Transformada}
  \[
    w^tDw+E^tw+f=0
  \]

  \[
    \begin{bmatrix}\mu&\nu\end{bmatrix}
    \mat{3&0\\0&-1}
    \mat{\mu\\\nu}
    +
    \begin{bmatrix}-2\sqrt{2}&0\end{bmatrix}
    \mat{\mu\\\nu}
    -3=0
  \]

  \[
    3\mu^2-\nu^2-2\sqrt{2}\mu-3=0
  \]
\end{resolution}

\begin{resolution}{Completando Quadrados}
  \[
    3\mu^2-\nu^2-2\sqrt{2}\mu-3=0
  \]
  \[
    (3\mu^2-2\sqrt{2}\mu)-\nu^2-3=0
  \]
  \[
    3\left(\mu^2-\frac{2\sqrt{2}}{3}\mu\right)
    -\nu^2-3=0
  \]

  \[
    \begin{aligned}
      \mu^2-\frac{2\sqrt{2}}{3}\mu
      &=
      \mu^2-\frac{2\sqrt{2}}{3}\mu
      +\left(\frac{\sqrt{2}}{3}\right)^2
      -\left(\frac{\sqrt{2}}{3}\right)^2\\
      &=
      \left(\mu-\frac{\sqrt{2}}{3}\right)^2-\frac{2}{9}
    \end{aligned}
  \]
\end{resolution}

\begin{resolution}{Completando Quadrados}
  \begin{columns}[T]
    \begin{column}{0.49\textwidth}
      \[
        3\left(\mu^2-\frac{2\sqrt{2}}{3}\mu\right)
        -\nu^2-3=0
      \]
      \[
        3\left[
        \left(\mu-\frac{\sqrt{2}}{3}\right)^2
        -\frac{2}{9}
        \right]-\nu^2=3
      \]
      \[
        3\left(\mu-\frac{\sqrt{2}}{3}\right)^2
        -\frac{2}{3}-\nu^2=3
      \]
      \[
        3\left(\mu-\frac{\sqrt{2}}{3}\right)^2
        -\nu^2=\frac{11}{3}
      \]
    \end{column}
    \begin{column}{0.49\textwidth}
      \[
        \frac{
        3\left(\mu-\frac{\sqrt{2}}{3}\right)^2
        }{11/3}
        -
        \frac{\nu^2}{11/3}=1
      \]
      \[
        \frac{
        \left(\mu-\frac{\sqrt{2}}{3}\right)^2
        }{11/9}
        -
        \frac{\nu^2}{11/3}=1
      \]
    \end{column}
  \end{columns}
\end{resolution}

\begin{resolution}{Identificação}
  \textbf{A cônica}
  \[
    \frac{\left(\mu-\frac{\sqrt{2}}{3}\right)^2}{11/9}
    -\frac{\nu^2}{11/3}=1
  \]

  é uma hipérbole com

  \textbf{Centro}
  \[
    C=\left(\frac{\sqrt{2}}{3},0\right)
    \approx(0{,}816;0)
  \]

  \textbf{Parâmetros}
  \[
    a=\sqrt{\frac{11}{9}}=\frac{\sqrt{11}}{3}
    \approx1{,}105
    \qquad
    b=\sqrt{\frac{11}{3}}
    \approx1{,}915
  \]
\end{resolution}


%------------------------------------------------------------------------------%
\begin{resolution}{Esboço}
  \centering
  \includegraphics{I-08-Rotacao-ex-2}
\end{resolution}
%------------------------------------------------------------------------------%

%==============================================================================
% Exemplo 2
%==============================================================================

%------------------------------------------------------------------------------%
\begin{question}
  Considerando a equação
  \[
    x^2 + 8xy + y^2 - 15 = 0
  \]
  \vspace{-\baselineskip}
  \begin{enumerate}
    \item escreva a equação na forma matricial
    \item diagonalize a matriz da parte quadrática
    \item faça a mudança de variáveis
    \item identifique a cônica
  \end{enumerate}
\end{question}

%------------------------------------------------------------------------------%
\begin{resolution}{Equação Matricial}
  Equação na forma matricial
  \[
    v^tAv+B^tv+f=0
  \]
  \[
    v=\begin{bmatrix}x\\y\end{bmatrix}
    \qquad
    A=\begin{bmatrix}
      1&4\\
      4&1
    \end{bmatrix}
  \]
  \[
    B=\begin{bmatrix}0\\0\end{bmatrix}
    \qquad
    f=-15
  \]
\end{resolution}

%------------------------------------------------------------------------------%
\begin{resolution}{Diagonalizando $A$ -- Polinômio Característico}
  \begin{align*}
    p(\lambda)
    & = \det(A-\lambda I)
    \\
    & = \begin{vmatrix}
      1-\lambda&4\\
      4&1-\lambda
    \end{vmatrix}
    \\
    & = (1-\lambda)^2-16
    \\
    & = \lambda^2-2\lambda-15
  \end{align*}
\end{resolution}

%------------------------------------------------------------------------------%
\begin{resolution}{Diagonalizando $A$ -- Autovalores}
  Equação característica
  \[
    \lambda^2-2\lambda-15=0
  \]
  \[
    (\lambda-5)(\lambda+3)=0
  \]
  Autovalores
  \[
    \lambda_1=5
    \qquad
    \lambda_2=-3
  \]
\end{resolution}

%------------------------------------------------------------------------------%
\begin{resolution}{Autovetores Associados a $\lambda_1=5$}
  \begin{align*}
    (A-5I)v_1 & = \vec{0} \\
    \begin{bmatrix}
      -4&4\\
      4&-4
    \end{bmatrix}
    \vetor{x}{y}
    & = \vetor{0}{0}
  \end{align*}
  Portanto
  \[
    -4x+4y=0
    \qquad
    y=x
  \]
  Escolhendo \(x=1\)
  \[
    v_1=\vetor{1}{1}
  \]
\end{resolution}

%------------------------------------------------------------------------------%
\begin{resolution}{Autovetores Associados a $\lambda_2=-3$}
  \begin{align*}
    (A+3I)v_2 & = \vec{0} \\
    \begin{bmatrix}
      4&4\\
      4&4
    \end{bmatrix}
    \vetor{x}{y}
    & = \vetor{0}{0}
  \end{align*}
  Portanto
  \[
    4x+4y=0
    \qquad
    y=-x
  \]
  Escolhendo \(x=1\)
  \[
    v_2=\vetor{1}{-1}
  \]
\end{resolution}

%------------------------------------------------------------------------------%
\begin{resolution}{Matriz Ortogonal}
  Normalizando os autovetores
  \[
    u_1=\frac{1}{\sqrt{2}}\vetor{1}{1}
    \qquad
    u_2=\frac{1}{\sqrt{2}}\vetor{1}{-1}
  \]
  Matriz ortogonal
  \[
    P=\frac{1}{\sqrt{2}}
    \begin{bmatrix}
      1&1\\
      1&-1
    \end{bmatrix}
  \]
  Matriz diagonal
  \[
    D=P^tAP
    =\begin{bmatrix}
      5&0\\
      0&-3
    \end{bmatrix}
  \]
\end{resolution}

%------------------------------------------------------------------------------%
\begin{resolution}{Mudança de Variáveis}
  Mudança de variáveis
  \[
    v=Pw
    \qquad
    w=\vetor{\mu}{\nu}
  \]
  \[
    \begin{bmatrix}
      x\\y
    \end{bmatrix}
    =
    \frac{1}{\sqrt{2}}
    \begin{bmatrix}
      1&1\\
      1&-1
    \end{bmatrix}
    \begin{bmatrix}
      \mu\\\nu
    \end{bmatrix}
  \]
  Portanto
  \[
    x=\frac{\mu+\nu}{\sqrt{2}}
    \qquad
    y=\frac{\mu-\nu}{\sqrt{2}}
  \]
\end{resolution}

%------------------------------------------------------------------------------%
\begin{resolution}{Equação Transformada}
  Como \(B=\vec{0}\), temos \(E=P^tB=\vec{0}\).

  \begin{align*}
    w^tDw+E^tw+f & = 0
    \\
    5\mu^2-3\nu^2-15 & = 0
    \\
    5\mu^2-3\nu^2 & = 15
  \end{align*}
\end{resolution}

%------------------------------------------------------------------------------%
\begin{resolution}{Identificação}
  Dividindo a equação por \(15\)
  \[
    \frac{\mu^2}{3}-\frac{\nu^2}{5}=1
  \]
  A cônica é uma hipérbole com centro na origem.

  Semieixos
  \[
    a=\sqrt{3}
    \qquad
    b=\sqrt{5}
  \]

  O eixo real está na direção de \(\mu\).
\end{resolution}

%------------------------------------------------------------------------------%
\begin{resolution}{Esboço}
  \centering
  \includegraphics{I-08-Rotacao-ex-3}
\end{resolution}
%------------------------------------------------------------------------------%


%------------------------------------------------------------------------------%
\begin{question}
  Considerando a equação
  \[
    x^2+2xy+y^2+2x-2y=0
  \]
  \vspace{-\baselineskip}
  \begin{enumerate}
    \item escreva a equação na forma matricial
    \item diagonalize a matriz da parte quadrática
    \item faça a mudança de variáveis
    \item identifique a cônica
  \end{enumerate}
\end{question}

%------------------------------------------------------------------------------%
\begin{resolution}{Equação Matricial}
  Equação na forma matricial
  \[
    v^tAv+B^tv+f=0
  \]
  \[
    v=\begin{bmatrix}x\\y\end{bmatrix}
    \qquad
    A=\begin{bmatrix}
      1&1\\
      1&1
    \end{bmatrix}
  \]
  \[
    B=\begin{bmatrix}2\\-2\end{bmatrix}
    \qquad
    f=0
  \]
\end{resolution}

%------------------------------------------------------------------------------%
\begin{resolution}{Diagonalizando $A$ -- Polinômio Característico}
  \begin{align*}
    p(\lambda)
    & = \det(A-\lambda I)
    \\
    & = \begin{vmatrix}
      1-\lambda&1\\
      1&1-\lambda
    \end{vmatrix}
    \\
    & = (1-\lambda)^2-1
    \\
    & = \lambda^2-2\lambda
  \end{align*}
\end{resolution}

%------------------------------------------------------------------------------%
\begin{resolution}{Diagonalizando $A$ -- Autovalores}
  Equação característica
  \[
    \lambda^2-2\lambda=0
  \]
  \[
    \lambda(\lambda-2)=0
  \]
  Autovalores
  \[
    \lambda_1=2
    \qquad
    \lambda_2=0
  \]
\end{resolution}

%------------------------------------------------------------------------------%
\begin{resolution}{Autovetores Associados a $\lambda_1=2$}
  \begin{align*}
    (A-2I)v_1 & = \vec{0} \\
    \begin{bmatrix}
      -1&1\\
      1&-1
    \end{bmatrix}
    \vetor{x}{y}
    & = \vetor{0}{0}
  \end{align*}
  Portanto
  \[
    -x+y=0
    \qquad
    y=x
  \]
  Escolhendo \(x=1\)
  \[
    v_1=\vetor{1}{1}
  \]
\end{resolution}

%------------------------------------------------------------------------------%
\begin{resolution}{Autovetores Associados a $\lambda_2=0$}
  \begin{align*}
    Av_2 & = \vec{0} \\
    \begin{bmatrix}
      1&1\\
      1&1
    \end{bmatrix}
    \vetor{x}{y}
    & = \vetor{0}{0}
  \end{align*}
  Portanto
  \[
    x+y=0
    \qquad
    y=-x
  \]
  Escolhendo \(x=1\)
  \[
    v_2=\vetor{1}{-1}
  \]
\end{resolution}

%------------------------------------------------------------------------------%
\begin{resolution}{Matriz Ortogonal}
  Normalizando os autovetores
  \[
    u_1=\frac{1}{\sqrt{2}}\vetor{1}{1}
    \qquad
    u_2=\frac{1}{\sqrt{2}}\vetor{1}{-1}
  \]
  Matriz ortogonal
  \[
    P=\frac{1}{\sqrt{2}}
    \begin{bmatrix}
      1&1\\
      1&-1
    \end{bmatrix}
  \]
  Matriz diagonal
  \[
    D=P^tAP
    =\begin{bmatrix}
      2&0\\
      0&0
    \end{bmatrix}
  \]
\end{resolution}

%------------------------------------------------------------------------------%
\begin{resolution}{Mudança de Variáveis}
  Mudança de variáveis
  \[
    v=Pw
    \qquad
    w=\vetor{\mu}{\nu}
  \]
  \[
    x=\frac{\mu+\nu}{\sqrt{2}}
    \qquad
    y=\frac{\mu-\nu}{\sqrt{2}}
  \]
\end{resolution}

%------------------------------------------------------------------------------%
\begin{resolution}{Transformando o Vetor $B$}
  \[
    E=P^tB
  \]
  \[
    E=
    \frac{1}{\sqrt{2}}
    \begin{bmatrix}
      1&1\\
      1&-1
    \end{bmatrix}
    \begin{bmatrix}
      2\\-2
    \end{bmatrix}
    =
    \begin{bmatrix}
      0\\2\sqrt{2}
    \end{bmatrix}
  \]
\end{resolution}

%------------------------------------------------------------------------------%
\begin{resolution}{Equação Transformada}
  \begin{align*}
    w^tDw+E^tw+f & = 0
    \\
    2\mu^2+2\sqrt{2}\nu & = 0
    \\
    \mu^2 & = -\sqrt{2}\nu
  \end{align*}
\end{resolution}

%------------------------------------------------------------------------------%
\begin{resolution}{Identificação}
  A equação
  \[
    \mu^2=-\sqrt{2}\nu
  \]
  representa uma parábola com vértice na origem.

  Seu eixo de simetria está na direção de \(\nu\).

  Como o coeficiente de \(\nu\) é negativo, a parábola
  tem concavidade voltada para o sentido negativo desse eixo.
\end{resolution}

%------------------------------------------------------------------------------%
\begin{resolution}{Esboço}
  \centering
  \includegraphics{I-08-Rotacao-ex-4}
\end{resolution}
%------------------------------------------------------------------------------%

%------------------------------------------------------------------------------%
\begin{question}
  Considerando a equação
  \[
    3x^2+4xy+3y^2+2x-2y-3=0
  \]
  \vspace{-\baselineskip}
  \begin{enumerate}
    \item escreva a equação na forma matricial
    \item diagonalize a matriz da parte quadrática
    \item faça a mudança de variáveis
    \item transforme a equação e complete os quadrados
    \item identifique a cônica
  \end{enumerate}
\end{question}

%------------------------------------------------------------------------------%
\begin{resolution}{Equação Matricial}
  Equação na forma matricial
  \[
    v^tAv+B^tv+f=0
  \]
  \[
    v=\begin{bmatrix}x\\y\end{bmatrix}
    \qquad
    A=\begin{bmatrix}
      3&2\\
      2&3
    \end{bmatrix}
  \]
  \[
    B=\begin{bmatrix}2\\-2\end{bmatrix}
    \qquad
    f=-3
  \]
\end{resolution}

%------------------------------------------------------------------------------%
\begin{resolution}{Diagonalizando $A$ -- Polinômio Característico}
  \begin{align*}
    p(\lambda)
    & = \det(A-\lambda I)
    \\
    & = \begin{vmatrix}
      3-\lambda&2\\
      2&3-\lambda
    \end{vmatrix}
    \\
    & = (3-\lambda)^2-4
    \\
    & = \lambda^2-6\lambda+5
  \end{align*}
\end{resolution}

%------------------------------------------------------------------------------%
\begin{resolution}{Diagonalizando $A$ -- Autovalores}
  Equação característica
  \[
    \lambda^2-6\lambda+5=0
  \]
  \[
    (\lambda-5)(\lambda-1)=0
  \]
  Autovalores
  \[
    \lambda_1=5
    \qquad
    \lambda_2=1
  \]
\end{resolution}

%------------------------------------------------------------------------------%
\begin{resolution}{Autovetores Associados a $\lambda_1=5$}
  \begin{align*}
    (A-5I)v_1 & = \vec{0} \\
    \begin{bmatrix}
      -2&2\\
      2&-2
    \end{bmatrix}
    \vetor{x}{y}
    & = \vetor{0}{0}
  \end{align*}
  \[
    -2x+2y=0
    \qquad
    y=x
  \]
  Escolhendo \(x=1\)
  \[
    v_1=\vetor{1}{1}
  \]
\end{resolution}

%------------------------------------------------------------------------------%
\begin{resolution}{Autovetores Associados a $\lambda_2=1$}
  \begin{align*}
    (A-I)v_2 & = \vec{0} \\
    \begin{bmatrix}
      2&2\\
      2&2
    \end{bmatrix}
    \vetor{x}{y}
    & = \vetor{0}{0}
  \end{align*}
  \[
    2x+2y=0
    \qquad
    y=-x
  \]
  Escolhendo \(x=1\)
  \[
    v_2=\vetor{1}{-1}
  \]
\end{resolution}

%------------------------------------------------------------------------------%
\begin{resolution}{Matriz Ortogonal}
  Normalizando os autovetores
  \[
    u_1=\frac{1}{\sqrt{2}}\vetor{1}{1}
    \qquad
    u_2=\frac{1}{\sqrt{2}}\vetor{1}{-1}
  \]
  Matriz ortogonal
  \[
    P=\frac{1}{\sqrt{2}}
    \begin{bmatrix}
      1&1\\
      1&-1
    \end{bmatrix}
  \]
  Matriz diagonal
  \[
    D=P^tAP
    =\begin{bmatrix}
      5&0\\
      0&1
    \end{bmatrix}
  \]
\end{resolution}

%------------------------------------------------------------------------------%
\begin{resolution}{Mudança de Variáveis}
  Mudança de variáveis
  \[
    v=Pw
    \qquad
    w=\vetor{\mu}{\nu}
  \]
  \[
    x=\frac{\mu+\nu}{\sqrt{2}}
    \qquad
    y=\frac{\mu-\nu}{\sqrt{2}}
  \]
\end{resolution}

%------------------------------------------------------------------------------%
\begin{resolution}{Transformando o Vetor $B$}
  \[
    E=P^tB
  \]
  \[
    E=
    \frac{1}{\sqrt{2}}
    \begin{bmatrix}
      1&1\\
      1&-1
    \end{bmatrix}
    \begin{bmatrix}
      2\\-2
    \end{bmatrix}
    =
    \begin{bmatrix}
      0\\2\sqrt{2}
    \end{bmatrix}
  \]
\end{resolution}

%------------------------------------------------------------------------------%
\begin{resolution}{Equação Transformada}
  \begin{align*}
    w^tDw+E^tw+f & = 0
    \\
    5\mu^2+\nu^2+2\sqrt{2}\nu-3 & = 0
  \end{align*}
\end{resolution}

%------------------------------------------------------------------------------%
\begin{resolution}{Completando Quadrados}
  Agrupando os termos
  \begin{align*}
    5\mu^2+\nu^2+2\sqrt{2}\nu-3 & = 0
    \\
    5\mu^2+\left(\nu^2+2\sqrt{2}\nu\right) & = 3
    \\
    5\mu^2+\left(\nu+\sqrt{2}\right)^2-2 & = 3
    \\
    5\mu^2+\left(\nu+\sqrt{2}\right)^2 & = 5
  \end{align*}
\end{resolution}

%------------------------------------------------------------------------------%
\begin{resolution}{Identificação}
  Dividindo a equação por \(5\)
  \[
    \mu^2+\frac{\left(\nu+\sqrt{2}\right)^2}{5}=1
  \]
  A cônica é uma elipse.

  Centro no sistema rotacionado
  \[
    C'=(0,-\sqrt{2})
  \]

  Centro no sistema original
  \[
    C=(-1,1)
  \]

  Semieixos
  \[
    a=1
    \qquad
    b=\sqrt{5}
  \]

  O eixo maior está na direção de \(\nu\).
\end{resolution}

%------------------------------------------------------------------------------%
\begin{resolution}{Esboço}
  \centering
  \includegraphics{I-08-Rotacao-ex-5}
\end{resolution}
%------------------------------------------------------------------------------%

%------------------------------------------------------------------------------%
\begin{question}
  Considerando a equação
  \[
    x^2+8xy+y^2+10x+10y-5=0
  \]
  \vspace{-\baselineskip}
  \begin{enumerate}
    \item escreva a equação na forma matricial
    \item diagonalize a matriz da parte quadrática
    \item faça a mudança de variáveis
    \item transforme a equação e complete os quadrados
    \item identifique a cônica
  \end{enumerate}
\end{question}

%------------------------------------------------------------------------------%
\begin{resolution}{Equação Matricial}
  Equação na forma matricial
  \[
    v^tAv+B^tv+f=0
  \]
  \[
    v=\begin{bmatrix}x\\y\end{bmatrix}
    \qquad
    A=\begin{bmatrix}
      1&4\\
      4&1
    \end{bmatrix}
  \]
  \[
    B=\begin{bmatrix}10\\10\end{bmatrix}
    \qquad
    f=-5
  \]
\end{resolution}

%------------------------------------------------------------------------------%
\begin{resolution}{Diagonalizando $A$ -- Polinômio Característico}
  \begin{align*}
    p(\lambda)
    & = \det(A-\lambda I)
    \\
    & = \begin{vmatrix}
      1-\lambda&4\\
      4&1-\lambda
    \end{vmatrix}
    \\
    & = (1-\lambda)^2-16
    \\
    & = \lambda^2-2\lambda-15
  \end{align*}
\end{resolution}

%------------------------------------------------------------------------------%
\begin{resolution}{Diagonalizando $A$ -- Autovalores}
  Equação característica
  \[
    \lambda^2-2\lambda-15=0
  \]
  \[
    (\lambda-5)(\lambda+3)=0
  \]
  Autovalores
  \[
    \lambda_1=5
    \qquad
    \lambda_2=-3
  \]
\end{resolution}

%------------------------------------------------------------------------------%
\begin{resolution}{Autovetores Associados a $\lambda_1=5$}
  \begin{align*}
    (A-5I)v_1 & = \vec{0} \\
    \begin{bmatrix}
      -4&4\\
      4&-4
    \end{bmatrix}
    \vetor{x}{y}
    & = \vetor{0}{0}
  \end{align*}
  \[
    -4x+4y=0
    \qquad
    y=x
  \]
  Escolhendo \(x=1\)
  \[
    v_1=\vetor{1}{1}
  \]
\end{resolution}

%------------------------------------------------------------------------------%
\begin{resolution}{Autovetores Associados a $\lambda_2=-3$}
  \begin{align*}
    (A+3I)v_2 & = \vec{0} \\
    \begin{bmatrix}
      4&4\\
      4&4
    \end{bmatrix}
    \vetor{x}{y}
    & = \vetor{0}{0}
  \end{align*}
  \[
    4x+4y=0
    \qquad
    y=-x
  \]
  Escolhendo \(x=1\)
  \[
    v_2=\vetor{1}{-1}
  \]
\end{resolution}

%------------------------------------------------------------------------------%
\begin{resolution}{Matriz Ortogonal}
  Normalizando os autovetores
  \[
    u_1=\frac{1}{\sqrt{2}}\vetor{1}{1}
    \qquad
    u_2=\frac{1}{\sqrt{2}}\vetor{1}{-1}
  \]
  Matriz ortogonal
  \[
    P=\frac{1}{\sqrt{2}}
    \begin{bmatrix}
      1&1\\
      1&-1
    \end{bmatrix}
  \]
  Matriz diagonal
  \[
    D=P^tAP
    =\begin{bmatrix}
      5&0\\
      0&-3
    \end{bmatrix}
  \]
\end{resolution}

%------------------------------------------------------------------------------%
\begin{resolution}{Mudança de Variáveis}
  Mudança de variáveis
  \[
    v=Pw
    \qquad
    w=\vetor{\mu}{\nu}
  \]
  \[
    x=\frac{\mu+\nu}{\sqrt{2}}
    \qquad
    y=\frac{\mu-\nu}{\sqrt{2}}
  \]
\end{resolution}

%------------------------------------------------------------------------------%
\begin{resolution}{Transformando o Vetor $B$}
  \[
    E=P^tB
  \]
  \[
    E=
    \frac{1}{\sqrt{2}}
    \begin{bmatrix}
      1&1\\
      1&-1
    \end{bmatrix}
    \begin{bmatrix}
      10\\10
    \end{bmatrix}
    =
    \begin{bmatrix}
      10\sqrt{2}\\0
    \end{bmatrix}
  \]
\end{resolution}

%------------------------------------------------------------------------------%
\begin{resolution}{Equação Transformada}
  \begin{align*}
    w^tDw+E^tw+f & = 0
    \\
    5\mu^2-3\nu^2+10\sqrt{2}\mu-5 & = 0
  \end{align*}
\end{resolution}

%------------------------------------------------------------------------------%
\begin{resolution}{Completando Quadrados}
  Agrupando os termos
  \begin{align*}
    5\mu^2+10\sqrt{2}\mu-3\nu^2-5 & = 0
    \\
    5\left(\mu^2+2\sqrt{2}\mu\right)-3\nu^2 & = 5
    \\
    5\left[\left(\mu+\sqrt{2}\right)^2-2\right]-3\nu^2 & = 5
    \\
    5\left(\mu+\sqrt{2}\right)^2-3\nu^2 & = 15
  \end{align*}
\end{resolution}

%------------------------------------------------------------------------------%
\begin{resolution}{Identificação}
  Dividindo a equação por \(15\)
  \[
    \frac{\left(\mu+\sqrt{2}\right)^2}{3}
    -\frac{\nu^2}{5}=1
  \]
  A cônica é uma hipérbole.

  Centro no sistema rotacionado
  \[
    C'=(-\sqrt{2},0)
  \]

  Centro no sistema original
  \[
    C=(-1,-1)
  \]

  Semieixos
  \[
    a=\sqrt{3}
    \qquad
    b=\sqrt{5}
  \]

  O eixo real está na direção de \(\mu\).
\end{resolution}


%------------------------------------------------------------------------------%
\begin{resolution}{Esboço}
  \centering
  \includegraphics{I-08-Rotacao-ex-6}
\end{resolution}
%------------------------------------------------------------------------------%

%------------------------------------------------------------------------------%
\begin{question}
  Considerando a equação
  \[
    x^2+2xy+y^2+4x-4y-8=0
  \]
  \vspace{-\baselineskip}
  \begin{enumerate}
    \item escreva a equação na forma matricial
    \item diagonalize a matriz da parte quadrática
    \item faça a mudança de variáveis
    \item transforme a equação
    \item identifique a cônica
  \end{enumerate}
\end{question}

%------------------------------------------------------------------------------%
\begin{resolution}{Equação Matricial}
  Equação na forma matricial
  \[
    v^tAv+B^tv+f=0
  \]
  \[
    v=\begin{bmatrix}x\\y\end{bmatrix}
    \qquad
    A=\begin{bmatrix}
      1&1\\
      1&1
    \end{bmatrix}
  \]
  \[
    B=\begin{bmatrix}4\\-4\end{bmatrix}
    \qquad
    f=-8
  \]
\end{resolution}

%------------------------------------------------------------------------------%
\begin{resolution}{Diagonalizando $A$ -- Polinômio Característico}
  \begin{align*}
    p(\lambda)
    & = \det(A-\lambda I)
    \\
    & = \begin{vmatrix}
      1-\lambda&1\\
      1&1-\lambda
    \end{vmatrix}
    \\
    & = (1-\lambda)^2-1
    \\
    & = \lambda^2-2\lambda
  \end{align*}
\end{resolution}

%------------------------------------------------------------------------------%
\begin{resolution}{Diagonalizando $A$ -- Autovalores}
  Equação característica
  \[
    \lambda^2-2\lambda=0
  \]
  \[
    \lambda(\lambda-2)=0
  \]
  Autovalores
  \[
    \lambda_1=2
    \qquad
    \lambda_2=0
  \]
\end{resolution}

%------------------------------------------------------------------------------%
\begin{resolution}{Autovetores Associados a $\lambda_1=2$}
  \begin{align*}
    (A-2I)v_1 & = \vec{0} \\
    \begin{bmatrix}
      -1&1\\
      1&-1
    \end{bmatrix}
    \vetor{x}{y}
    & = \vetor{0}{0}
  \end{align*}
  \[
    -x+y=0
    \qquad
    y=x
  \]
  Escolhendo \(x=1\)
  \[
    v_1=\vetor{1}{1}
  \]
\end{resolution}

%------------------------------------------------------------------------------%
\begin{resolution}{Autovetores Associados a $\lambda_2=0$}
  \begin{align*}
    Av_2 & = \vec{0} \\
    \begin{bmatrix}
      1&1\\
      1&1
    \end{bmatrix}
    \vetor{x}{y}
    & = \vetor{0}{0}
  \end{align*}
  \[
    x+y=0
    \qquad
    y=-x
  \]
  Escolhendo \(x=1\)
  \[
    v_2=\vetor{1}{-1}
  \]
\end{resolution}

%------------------------------------------------------------------------------%
\begin{resolution}{Matriz Ortogonal}
  Normalizando os autovetores
  \[
    u_1=\frac{1}{\sqrt{2}}\vetor{1}{1}
    \qquad
    u_2=\frac{1}{\sqrt{2}}\vetor{1}{-1}
  \]
  Matriz ortogonal
  \[
    P=\frac{1}{\sqrt{2}}
    \begin{bmatrix}
      1&1\\
      1&-1
    \end{bmatrix}
  \]
  Matriz diagonal
  \[
    D=P^tAP
    =\begin{bmatrix}
      2&0\\
      0&0
    \end{bmatrix}
  \]
\end{resolution}

%------------------------------------------------------------------------------%
\begin{resolution}{Mudança de Variáveis}
  Mudança de variáveis
  \[
    v=Pw
    \qquad
    w=\vetor{\mu}{\nu}
  \]
  \[
    x=\frac{\mu+\nu}{\sqrt{2}}
    \qquad
    y=\frac{\mu-\nu}{\sqrt{2}}
  \]
\end{resolution}

%------------------------------------------------------------------------------%
\begin{resolution}{Transformando o Vetor $B$}
  \[
    E=P^tB
  \]
  \[
    E=
    \frac{1}{\sqrt{2}}
    \begin{bmatrix}
      1&1\\
      1&-1
    \end{bmatrix}
    \begin{bmatrix}
      4\\-4
    \end{bmatrix}
    =
    \begin{bmatrix}
      0\\4\sqrt{2}
    \end{bmatrix}
  \]
\end{resolution}

%------------------------------------------------------------------------------%
\begin{resolution}{Equação Transformada}
  \begin{align*}
    w^tDw+E^tw+f & = 0
    \\
    2\mu^2+4\sqrt{2}\nu-8 & = 0
    \\
    \mu^2+2\sqrt{2}\nu & = 4
  \end{align*}
\end{resolution}

%------------------------------------------------------------------------------%
\begin{resolution}{Identificação}
  Isolando o termo quadrático
  \[
    \mu^2=-2\sqrt{2}\left(\nu-\sqrt{2}\right)
  \]
  A cônica é uma parábola.

  Vértice no sistema rotacionado
  \[
    V'=(0,\sqrt{2})
  \]

  Vértice no sistema original
  \[
    V=(1,-1)
  \]

  O eixo de simetria está na direção de \(\nu\).

  A parábola tem concavidade voltada para o sentido negativo
  desse eixo.
\end{resolution}

%------------------------------------------------------------------------------%
\begin{resolution}{Esboço}
  \centering
  \includegraphics{I-08-Rotacao-ex-7}
\end{resolution}
%------------------------------------------------------------------------------%


%------------------------------------------------------------------------------%
%\section{Lista Mínima}
%
%%------------------------------------------------------------------------------%
%\begin{frame}{Lista Mínima}
%  \begin{enumerate}
%    \item
%  \end{enumerate}
%
%  \vfill
%  Atenção: A prova é baseada no livro, não nas apresentações
%\end{frame}

%------------------------------------------------------------------------------%
\end{document}
%------------------------------------------------------------------------------%
